\documentclass[a4paper,11pt]{article}
\usepackage[utf8]{inputenc}
\usepackage[T1]{fontenc}
\usepackage[italian]{babel}
\usepackage{geometry}
\usepackage{hyperref}
\usepackage{xcolor}
\usepackage{enumitem}
\usepackage{booktabs}
\usepackage[protrusion=true,expansion=false]{microtype}

\geometry{margin=2.7cm}
\setlength{\parskip}{0.5em}
\setlength{\parindent}{0pt}

\newenvironment{citazione}{\begin{quote}\itshape}{\end{quote}}

\hypersetup{colorlinks=true, linkcolor=blue!50!black, urlcolor=blue!50!black}

\title{\textbf{The Atlas of AI} e \textbf{Le Macchine Sapienti}\\[0.3em]
\large Kate Crawford (2021) e Paolo Benanti -- due letture a confronto\\[0.3em]
\normalsize Collegamenti con il corso \emph{Etica, Società e Privacy}\\
Università di Torino -- Magistrale Informatica -- A.A. 2024/2025}
\author{}
\date{}

\begin{document}
\maketitle
\tableofcontents
\newpage

%=============================================================================
\section{Due autori, due punti di partenza}
%=============================================================================

Questa relazione raccoglie due libri che arrivano alla stessa domanda -- cosa sta davvero facendo l'intelligenza artificiale alla nostra società -- partendo da due punti di vista quasi opposti, e che per questo vale la pena leggere insieme piuttosto che separatamente.

Kate Crawford è ricercatrice tra Microsoft Research, la New York University e la University of Southern California, e si occupa da oltre un decennio delle implicazioni sociali, politiche e ambientali dei sistemi di intelligenza artificiale. \emph{The Atlas of AI} (Yale University Press, 2021) nasce da anni di ricerca sul campo -- miniere di litio, data center, archivi craniometrici, basi di lancio spaziale -- e il suo metodo è prima di tutto materiale e politico: per Crawford l'IA va capita seguendo la materia di cui è fatta (minerali, energia, corpi di lavoratori, dati estratti senza consenso) e il potere che quella materia rende possibile.

Paolo Benanti è un frate francescano, teologo morale e bioeticista, docente alla Pontificia Università Gregoriana, consulente di papa Francesco e del governo italiano sui temi dell'intelligenza artificiale, e inventore del termine \emph{algoretica} per indicare un'etica degli algoritmi. \emph{Le Macchine Sapienti} affronta lo stesso terreno -- lavoro, sanità, giustizia, sorveglianza, guerra -- ma con un metodo diverso, che parte dalla domanda su cosa renda l'essere umano umano, e arriva a proporre codici etici e forme di governance capaci di tenere le macchine subordinate all'uomo.

È un confronto produttivo proprio perché i due autori non sono d'accordo su tutto: Crawford diffida delle carte etiche e dei codici di condotta, e propone di guardare al potere più che alla morale; Benanti scommette invece sulla possibilità di tradurre principi etici in regole tecniche e in governance pubblica. Leggerli insieme permette di vedere lo stesso problema -- l'illusione che l'IA sia neutra, oggettiva, immateriale -- da due angolazioni diverse: una più politica ed economica, l'altra più filosofica e teologica.

%=============================================================================
\section{Il libro di Crawford, capitolo per capitolo}
%=============================================================================

Il libro si apre con un aneddoto che vale la pena ricordare per intero, perché ne anticipa la tesi. A Berlino, nei primi anni del Novecento, un cavallo chiamato Hans diventò famoso perché sembrava in grado di fare calcoli, battendo lo zoccolo il numero esatto di volte corrispondente alla soluzione di un'operazione. Un'indagine più attenta rivelò che Hans non sapeva contare: leggeva segnali corporei impercettibili e involontari del suo addestratore. Crawford usa questa storia come metafora di un rischio che corriamo ancora oggi: scambiare per ``intelligenza'' un sistema che in realtà risponde a segnali nascosti nell'ambiente che lo circonda -- nel caso dell'IA contemporanea, enormi infrastrutture di estrazione mineraria, energia, lavoro umano e dati personali che restano invisibili a chi usa un assistente vocale o un sistema di riconoscimento facciale.

Da qui nasce la tesi che attraversa tutto il libro: l'intelligenza artificiale non è né artificiale né intelligente. Non è artificiale perché è fatta di materia -- minerali rari, energia, acqua, corpi di lavoratori. Non è intelligente perché i sistemi che chiamiamo così sono inferenza statistica addestrata su grandi quantità di dati, non pensiero. L'idea-guida che lega tutto il libro, organizzato come un atlante di luoghi e immagini parziali, è che l'IA sia un registro di potere: non crea potere dal nulla, ma registra, amplifica e rende invisibili rapporti di forza già esistenti, presentandoli come il semplice output neutro di un calcolo.

\subsection{Earth: la terra che l'IA consuma}

Il primo capitolo mostra come l'intelligenza artificiale, venduta come qualcosa di immateriale che vive nel cloud, dipenda in realtà da un'industria estrattiva fatta della stessa materia di quella mineraria e petrolifera. Crawford visita Silver Peak, in Nevada, l'unica miniera di litio ancora attiva negli Stati Uniti, e mostra il contrasto tra il paesaggio desertico devastato dall'estrazione e l'immagine pulita con cui Tesla vende le proprie batterie. Il racconto si sposta poi nella Repubblica Democratica del Congo e in Indonesia, dove l'estrazione di cobalto e stagno -- materiali indispensabili per smartphone e server -- è legata a violenza armata e ai cosiddetti ``minerali di conflitto''. A questa geografia mineraria si aggiunge quella dei data center: la base della NSA a Bluffdale, nello Utah, che consuma quantità enormi d'acqua in pieno deserto, e l'impianto Google a The Dalles, in Oregon, dipendente da infrastrutture pubbliche sussidiate che restano nascoste dietro la metafora del ``cloud''. Il filo che lega tutti questi casi è il mito della tecnologia pulita: l'illusione che il digitale sia leggero e immateriale, quando in realtà -- per usare le parole di Crawford -- l'IA ``nasce da laghi salati in Bolivia e miniere in Congo''.

\subsection{Labor: il lavoro nascosto dietro l'automazione}

Il secondo capitolo smonta l'illusione secondo cui i sistemi di IA funzionerebbero in modo automatico, senza bisogno di mani umane. Crawford descrive i magazzini di Amazon, dove i lavoratori sono tracciati da scanner portatili e i ritmi sono dettati da algoritmi in una versione digitale del taylorismo, e racconta lo sciopero organizzato nel centro logistico di Eagan, in Minnesota, da una comunità di lavoratori somalo-americani guidati da Abdi Muse e dall'Awood Center. Il caso più rilevante è quello di Amazon Mechanical Turk, la piattaforma che scompone compiti complessi in microattività pagate pochi centesimi -- etichettare immagini, risolvere CAPTCHA, moderare contenuti -- svolte da milioni di persone che restano invisibili all'utente finale. Crawford riprende da Mary Gray e Sid Suri il termine \emph{ghost work} e da Astra Taylor il termine \emph{fauxtomation} per indicare quanto spesso ciò che viene venduto come ``automatico'' nasconda lavoro umano dietro le quinte.

\subsection{Data: dataset che non sono mai neutri}

Il terzo capitolo mostra che i dataset non sono mai semplice materiale grezzo, ma costruzioni storiche e politiche, spesso raccolte senza consenso. Crawford ricostruisce la genealogia di questa pratica a partire da Francis Galton e Alphonse Bertillon, pionieri ottocenteschi della classificazione biometrica e criminologica. Il caso più importante è quello di ImageNet, il dataset creato nel 2009 da Fei-Fei Li raschiando milioni di immagini dai motori di ricerca ed etichettandole tramite Amazon Mechanical Turk: una dimostrazione di come estrazione di lavoro ed estrazione di dati siano in realtà la stessa cosa vista da due angolazioni diverse. Crawford racconta anche dataset come DukeMTMC o MS-Celeb, costruiti raccogliendo foto pubbliche senza permesso e poi ritirati dopo scandali di privacy, per parlare di quello che chiama ``la fine del consenso''.

\subsection{Classification: classificare è un atto di potere}

Il quarto capitolo si apre con un'immagine forte: Crawford visita al Penn Museum di Philadelphia la collezione di circa mille teschi raccolti dal craniologo Samuel Morton, morto nel 1851, che riempiva le cavità craniche con pallini di piombo per ``dimostrare'' la superiorità razziale bianca. La classificazione algoritmica contemporanea, sostiene Crawford, non è semplicemente simile a queste pratiche pseudoscientifiche, ma ne è in molti casi una continuazione diretta sotto altra forma tecnica. Il capitolo lo dimostra ripercorrendo le categorie problematiche di ImageNet, rese visibili dal progetto artistico \emph{ImageNet Roulette} di Trevor Paglen, e lo studio \emph{Gender Shades} di Joy Buolamwini e Timnit Gebru, che nel 2018 dimostrò tassi di errore molto più alti nel riconoscimento facciale di donne con pelle scura. Il caso dello strumento di selezione del personale sperimentato da Amazon, che continuava a penalizzare le candidate donne anche dopo la rimozione del riferimento esplicito al genere, è la prova, per Crawford, che il pregiudizio non è un difetto isolabile, ma una caratteristica intrinseca dell'atto stesso di classificare.

\subsection{Affect: leggere le emozioni dai volti}

Il quinto capitolo racconta la storia fragile, sul piano scientifico, dell'industria del riconoscimento delle emozioni. Tutto parte dal lavoro dello psicologo Paul Ekman, che nel 1967 si recò in Papua Nuova Guinea per dimostrare l'universalità delle espressioni facciali delle emozioni, con una ricerca finanziata dall'agenzia militare ARPA. Da quel lavoro nasce il Facial Action Coding System, applicato oggi da aziende come HireVue o Affectiva. Una revisione del 2019 guidata da Lisa Feldman Barrett conclude però che non è possibile inferire con sicurezza la felicità da un sorriso o la rabbia da un'espressione corrucciata: l'industria, osserva Crawford, ha scelto la teoria che si adattava agli strumenti disponibili, non quella più solida.

\subsection{State: l'intelligenza artificiale e il potere statale}

Il sesto capitolo segue il filo che collega la sorveglianza militare a quella della vita quotidiana. Crawford riprende l'archivio Snowden per raccontare programmi della NSA come TREASUREMAP, e il caso di Project Maven, il programma del Pentagono per analizzare i filmati dei drone tramite l'IA di Google, che provocò una protesta interna firmata da oltre tremila dipendenti. Racconta poi Palantir, da contractor dei servizi di intelligence a fornitrice di sorveglianza per l'agenzia americana per l'immigrazione e per diversi dipartimenti di polizia locali. Il filo che lega questi episodi è l'infiltrazione dal basso degli strumenti nati per l'intelligence militare: tecnologie pensate per i campi di battaglia finiscono, con sorprendente rapidità, per essere usate nelle scuole, nelle stazioni di polizia, negli uffici per la disoccupazione.

\subsection{Power e Space: la conclusione e la coda}

Nelle pagine conclusive Crawford introduce il concetto di \emph{enchanted determinism}, coniato con lo storico Alex Campolo a partire dal caso di AlphaGo Zero: l'industria presenta i propri sistemi con un doppio registro retorico, insieme incantato (oltre la comprensione umana) e deterministico (capace di prevedere con certezza), un doppio registro che mistifica il potere reale dei sistemi di IA e li sottrae alla critica. Crawford critica come insufficienti le carte etiche pubblicate da aziende e governi, e propone di spostare l'attenzione dall'etica al potere, rivendicando un diritto al rifiuto -- la possibilità di non usare un certo sistema in un certo contesto -- e citando la frase di Audre Lorde secondo cui gli strumenti del padrone non smantelleranno mai la casa del padrone. Il libro si chiude con la coda dedicata alla corsa allo spazio di Bezos e Musk, presentata come l'estensione finale della stessa logica estrattiva raccontata nel primo capitolo.

%=============================================================================
\section{Il libro di Benanti, capitolo per capitolo}
%=============================================================================

\emph{Le Macchine Sapienti} si apre con un paragone storico: Benanti accosta la diffusione dell'intelligenza artificiale alla scoperta di Galileo, ignorata a lungo dai filosofi aristotelici perché metteva in crisi un intero modo di percepire il mondo. Oggi, sostiene l'autore, ci troviamo in una situazione simile: lo sviluppo dell'IA non cambia solo gli strumenti che usiamo, ma la nostra stessa percezione e cognizione della realtà.

\subsection{Capitolo 1: i nuovi terreni dell'intelligenza artificiale}

Benanti passa in rassegna gli ambiti in cui l'IA sta già cambiando la vita quotidiana. Nel mondo del \textbf{lavoro}, le macchine capaci di decisioni autonome aprono per la prima volta una vera competizione con l'uomo nel processo produttivo, non solo nella forza fisica ma nella decisione stessa. Nella \textbf{sanità}, applicazioni come ginger.io raccolgono dati di movimento, chiamate e messaggi per segnalare ai medici cambiamenti nel comportamento di pazienti cronici -- un uso promettente dei dati personali, ma che pone subito la domanda di chi decida cosa sia ``normale''. Sulle \textbf{tecnologie riproduttive}, Benanti racconta sistemi che usano dati fotografici o livelli di DNA per selezionare quali embrioni impiantare, e chiede esplicitamente se si possa affidare a una macchina la decisione su chi sia degno di nascere.

Il capitolo si sofferma poi con particolare attenzione sulla \textbf{giustizia predittiva}: algoritmi che analizzano sentenze passate e profili dei condannati per calcolare il rischio di recidiva. L'argomento a favore di questi sistemi sembra impeccabile -- i giudici umani sono stanchi, hanno pregiudizi inconsci, sono influenzati da fattori esterni -- ma Benanti individua tre problemi precisi. Il primo è quello che chiama ``il passato che imprigiona il futuro'': se i dati storici contengono pregiudizi sociali, l'algoritmo li trasforma in una regola scientifica apparente, automatizzando e amplificando il pregiudizio umano invece di eliminarlo. Il secondo è la scatola nera: una sentenza giusta deve essere motivata, ma il calcolo che produce un punteggio di rischio resta incomprensibile sia al giudice sia all'imputato. Il terzo è la perdita dell'unicità della persona: l'imputato non viene più giudicato per quello che ha fatto, ma trattato come un elemento di una categoria statistica.

Da qui Benanti passa al caso, per lui esemplare, dei \textbf{sistemi di credito sociale} cinesi, nati nel 2014 da un documento del Consiglio di Stato che proponeva di misurare la fiducia nelle persone con un punteggio unico, il \emph{Citizen Score}, costruito monitorando acquisti, spostamenti, tempo passato online, bollette pagate -- e influenzato persino dal comportamento delle persone con cui ci si circonda. Il capitolo prosegue con un'analisi dei \textbf{social network} e di quello che Benanti chiama un mondo post-fattuale: l'economia dell'attenzione spinge le piattaforme a mostrare i contenuti che generano le reazioni più forti -- indignazione, rabbia, paura -- mentre l'algoritmo di personalizzazione costruisce attorno a ciascuno di noi una bolla, una \emph{echo chamber}, che isola dal confronto con opinioni diverse e radicalizza le posizioni. La crisi è aggravata dalla perdita dell'intermediazione: il post di un esperto e la teoria del complotto di un utente qualunque hanno oggi la stessa dignità visiva su una bacheca.

Il capitolo si chiude su due fronti distinti. Il primo è quello di \textbf{sicurezza nazionale e guerra}: droni e sistemi di fuoco robotici sempre più autonomi pongono il problema di chi sia responsabile quando una macchina, definita un \emph{Artificial Moral Agent}, prende decisioni che causano vittime civili. Il secondo è una riflessione sui ``nuovi scenari'': Benanti descrive una ``grande mediazione'' dell'esperienza umana, in cui ogni scelta -- un film, un lavoro, una relazione -- passa attraverso un algoritmo che decide cosa sia meglio per noi, e un ``capitalismo di sorveglianza'' fondato sulla raccolta sistematica di dati, che crea un'asimmetria di potere enorme tra le poche aziende che sanno tutto di noi e gli utenti che non sanno nulla del funzionamento dei loro sistemi.

\subsection{Capitolo 2: la condizione tecno-umana}

Nel secondo capitolo Benanti smonta l'idea che la tecnologia sia un semplice insieme di strumenti neutri che l'uomo usa a piacimento. Per l'autore l'essere umano è costitutivamente tecnologico: biologicamente fragile e incompleto -- senza gli artigli del leone, la velocità del ghepardo, la pelliccia dell'orso -- la specie umana ha dovuto creare tecnologia (la lancia, il fuoco, l'agricoltura) per sopravvivere. La tecnologia non è quindi qualcosa che si aggiunge alla vita umana, ma l'ambiente in cui l'uomo si realizza. Quello che cambia con l'intelligenza artificiale, sostiene Benanti, è che per la prima volta la tecnologia non amplifica più un'azione fisica (la ruota potenzia le gambe, la gru le braccia), ma si sostituisce al pensiero e alla scelta: la macchina elabora informazioni, impara dagli errori, decide.

Questa svolta mette in crisi la definizione classica di essere umano come unico animale dotato di ragione, e qui Benanti introduce una distinzione che è il cuore filosofico del libro: la differenza tra \textbf{sintassi e semantica}. Le macchine lavorano sulla sintassi -- collegano simboli e parole secondo regole statistiche, calcolando quale parola ha più probabilità di venire dopo l'altra, senza sapere cosa stiano dicendo. Gli esseri umani vivono invece di semantica: comprendono il significato, il peso emotivo ed etico delle parole e delle azioni. Una macchina può simulare la tristezza scrivendo un testo malinconico, ma non sa cosa sia il dolore. A questa distinzione Benanti aggiunge il tema del \textbf{corpo}: siamo fatti di carne, sangue, emozioni, e soprattutto siamo mortali; la nostra intelligenza non è separata dal corpo, passa attraverso la fatica, la fame, la paura di morire. Una macchina senza corpo non sperimenta vulnerabilità, e la sua intelligenza resta per questo fredda e priva della compassione che nasce dal riconoscersi simili nella fragilità. Da qui discende anche la \textbf{responsabilità etica}: solo l'uomo può decidere cosa sia giusto e sbagliato assumendosene la responsabilità, anche contro la logica dell'utilità; la macchina ottimizza un risultato, non ha coscienza morale, e proprio i limiti e gli errori umani sono spesso la fonte della creatività e della capacità di perdonare.

Il capitolo si chiude con due interludi che usano metafore efficaci. Il primo è quello della mappa: una mappa, per essere utile, deve semplificare la realtà, e i dati che alimentano l'IA sono soltanto una mappa della società, non la società stessa -- il rischio, scrive Benanti, è di confondere la mappa con il territorio, giudicando le persone solo in base ai loro dati passati ed escludendo l'imprevisto e il cambiamento. Il secondo interludio riguarda le ``macchine emotivo-razionali'', sistemi che analizzano tono di voce, espressioni facciali e battito cardiaco per dedurre stati emotivi: anche qui si tratta di un calcolo statistico che associa un'espressione a un'etichetta, non di una vera comprensione, ma il rischio di manipolazione commerciale e psicologica che ne deriva è reale, tanto più se queste macchine arrivano a sostituire relazioni umane complesse con surrogati sempre accoglienti e programmati per compiacere.

\subsection{Capitolo 3: un codice etico per le intelligenze artificiali}

Nel terzo capitolo Benanti propone un paradigma di \textbf{cooperazione}: l'obiettivo dell'IA deve essere aumentare la capacità cognitiva dell'uomo, non sostituirsi a essa -- la macchina supporta, l'uomo decide. Per realizzare questa cooperazione propone una ``moralità programmata'' fondata su tre principi: non ledere (l'IA non deve recare danno agli esseri umani), essere relazionale (la macchina resta subordinata all'uomo) e garantire la centralità umana (l'uomo non deve ridursi a un esecutore dei comandi della macchina). Riprendendo studi neuroscientifici, Benanti introduce poi il concetto di \emph{Free Won't}: la libertà umana si manifesta spesso come capacità di porre un veto a un impulso automatico, e il codice etico dovrebbe permettere a un uomo o a un sistema di interrompere l'automatismo algoritmico in nome di fini morali che la macchina, vincolata al solo calcolo ottimizzato, non può conoscere. Il capitolo elenca infine quattro requisiti tecnici che il codice etico dovrebbe imporre: intelligibilità (rifiuto della logica a scatola nera), intuizione (capacità di interfacciarsi con modalità conoscitive non solo logico-matematiche), adattabilità (capacità di adeguarsi al contesto reale) e regolazione (parametri di azione definiti da norme chiare).

\subsection{Capitolo 4: verso una governance delle intelligenze artificiali}

Se l'algoretica -- termine che Benanti ha contribuito a diffondere -- è la teoria che inserisce valori nel codice, la \textbf{governance} ne è la pratica politica e giuridica. Benanti sostiene che l'IA non è una tecnologia neutrale, ma una forma di potere capace di strutturare la società, e che per questo non basta l'autoregolamentazione delle aziende: serve una governance pubblica, globale e vincolante, capace di evitare che gli algoritmi sostituiscano le leggi dello stato e che le multinazionali tecnologiche diventino i veri sovrani del mondo contemporaneo. Parla esplicitamente di un rischio di \textbf{asimmetria di potere} e di \textbf{sovranità tecnologica}: quando un algoritmo decide chi ha accesso a un servizio o come si forma l'opinione pubblica, sta esercitando una funzione che prima apparteneva solo allo stato, e la governance serve a riaffermare il primato della politica sull'efficienza tecnica. Tra gli strumenti concreti propone un approccio basato sul rischio -- vincoli più severi solo dove il pericolo per i diritti fondamentali è più alto, come per le armi autonome o la classificazione sociale -- e la certificazione e l'audit indipendente dei sistemi prima della loro immissione sul mercato. Guarda infine con favore ai tentativi di regolamentazione internazionale, dalle linee guida dell'ONU e dell'OCSE all'approccio dell'Unione Europea, come tentativo di costruire un nucleo minimo di diritti condivisi tra culture diverse.

%=============================================================================
\section{Temi che attraversano entrambi i libri}
%=============================================================================

Messi a confronto, i due libri non si limitano a parlare delle stesse cose: si illuminano a vicenda, e in alcuni punti si contraddicono in modo produttivo.

\subsection{Il corpo: due usi opposti della stessa categoria}

È forse il punto di contatto più sorprendente. Crawford insiste sul corpo per smascherare l'IA: dietro la presunta immaterialità del cloud ci sono sempre corpi -- i minatori del Congo, i lavoratori di Amazon Mechanical Turk, i teschi di Samuel Morton -- e il suo intero progetto consiste nel rendere visibile questa materialità nascosta. Benanti insiste sul corpo per delimitare l'umano: ciò che rende un essere umano diverso da una macchina, scrive, è proprio l'avere un corpo fatto di fame, fatica, desiderio e paura della morte, qualcosa che nessuna intelligenza artificiale potrà mai possedere. Le due prospettive non sono in contraddizione, ma si completano: Crawford mostra che l'IA \emph{ha} un corpo che il marketing tecnologico nasconde (le miniere, i magazzini, i data center); Benanti mostra che l'IA \emph{non ha} il corpo che conta davvero per essere umani (la fragilità, la mortalità, il dolore vissuto in prima persona). Tenute insieme, le due tesi smontano da entrambi i lati la stessa illusione: quella di un'intelligenza artificiale che fluttua come pura mente, senza materia e senza limiti.

\subsection{Sintassi e semantica, ovvero la stessa critica con parole diverse}

La distinzione di Benanti tra sintassi (ciò che la macchina calcola) e semantica (ciò che solo l'uomo comprende) è, da una prospettiva filosofica e teologica, lo stesso argomento con cui Crawford apre il proprio libro affermando che l'IA non è intelligente, ma inferenza statistica priva di comprensione. È anche, detto con altre parole, l'argomento dei cosiddetti ``pappagalli stocastici'': i language model non sanno cosa dicono, ripetono sequenze di parole statisticamente probabili senza comprenderne il significato. Le due tradizioni di pensiero -- l'antropologia teologica di Benanti e la critica materialista di Crawford -- arrivano così, per strade indipendenti, alla medesima diagnosi: scambiamo per pensiero ciò che è soltanto un calcolo molto sofisticato di probabilità.

\subsection{Il passato che imprigiona il futuro: classificazione e giustizia predittiva}

Il capitolo di Crawford sulla classificazione e il capitolo di Benanti sulla giustizia predittiva raccontano, con esempi diversi, lo stesso meccanismo. Crawford lo mostra attraverso la genealogia che lega la craniometria di Samuel Morton alle categorie discriminatorie di ImageNet: un pregiudizio storico, una volta incorporato in un sistema tecnico, smette di apparire come pregiudizio e diventa una verità apparentemente oggettiva. Benanti descrive esattamente questo stesso processo applicandolo alla giustizia penale: se il passato contiene pregiudizi sociali, l'algoritmo che impara da quel passato trasforma il pregiudizio in una ``regola scientifica'', amplificandolo invece di correggerlo. Anche il sistema di credito sociale cinese descritto da Benanti e i casi di sorveglianza statale raccontati da Crawford nel capitolo sullo Stato condividono la stessa struttura: un punteggio che riduce una persona a un numero, costruito su dati storici che quella stessa persona non può contestare né correggere.

\subsection{Sorveglianza, potere e sovranità tecnologica}

Il ``capitalismo di sorveglianza'' di cui parla Benanti e il ``registro di potere'' che per Crawford definisce l'intelligenza artificiale descrivono lo stesso fenomeno da due lessici diversi: pochissime aziende private accumulano una conoscenza sterminata sui comportamenti, le paure e i desideri di miliardi di persone, mentre quelle stesse persone non sanno quasi nulla del funzionamento dei sistemi che le osservano. Quando Benanti parla di un vuoto di potere statale colmato dalle piattaforme tecnologiche, e quando Crawford racconta come tecnologie nate per l'intelligence militare finiscano per governare scuole e uffici di disoccupazione, stanno descrivendo la medesima erosione del confine tra potere pubblico e potere privato.

\subsection{Guerra autonoma e la domanda sulla responsabilità}

Il concetto di \emph{Artificial Moral Agent} introdotto da Benanti -- e la domanda su chi sia responsabile quando un sistema autonomo causa vittime civili -- trova un riscontro diretto nel capitolo di Crawford sullo Stato, dove si racconta come i drone strike vengano spesso decisi sulla base di metadati e pattern comportamentali piuttosto che di un'identificazione certa, fino alla frase, citata da Crawford, di un generale americano: ``uccidiamo persone sulla base di metadati''. Dove Crawford documenta il fatto, Benanti pone la domanda etica che quel fatto lascia aperta: una macchina che decide può essere moralmente responsabile, oppure la responsabilità resta sempre e comunque umana?

\subsection{Il punto di divergenza più importante: algoretica e governance contro politica del rifiuto}

È qui che i due libri si separano più nettamente, e proprio per questo vale la pena discuterne esplicitamente. Benanti crede che si possa costruire una ``moralità programmata'', incorporare principi etici nel codice e affidarsi a una governance pubblica -- audit, certificazioni, approccio basato sul rischio, accordi internazionali -- per tenere le macchine subordinate all'uomo. Crawford è molto più scettica verso questo tipo di soluzioni: nella conclusione del suo libro osserva che esistono già più di un centinaio di carte etiche sull'intelligenza artificiale pubblicate da aziende e governi, quasi tutte prive di qualsiasi effetto vincolante, e per questo propone di spostare l'attenzione dall'etica al potere, fino a rivendicare un diritto al rifiuto -- la possibilità per una comunità di non usare affatto un certo sistema. Non è una contraddizione che vada risolta a favore dell'uno o dell'altro: è piuttosto la prova che il problema dell'intelligenza artificiale ha sia una dimensione tecnica e normativa, su cui l'approccio di Benanti offre strumenti concreti, sia una dimensione di potere economico e politico, su cui l'approccio di Crawford resta insostituibile.

%=============================================================================
\section{Collegamenti con il corso di Etica}
%=============================================================================

\subsection{Contro l'inevitabilismo tecnologico}

Il corso critica la retorica dell'inevitabilismo tecnologico, quella secondo cui l'invasività digitale sarebbe una conseguenza naturale del progresso. Crawford arriva quasi alle stesse parole nelle pagine conclusive del libro, quando scrive che bisogna opporsi alle narrazioni dell'inevitabilità tecnologica secondo cui ``se qualcosa può essere fatto, sarà fatto''. Benanti, da una prospettiva diversa, dice in fondo la stessa cosa quando insiste che l'IA non è una tecnologia neutrale ma una forma di potere che richiede scelte politiche consapevoli: in entrambi i casi, il punto è che il futuro dell'intelligenza artificiale non è scritto in anticipo, ma è il risultato di decisioni che si possono ancora discutere e correggere.

\subsection{Una lista degli appunti che riassume Earth e Labor}

Tra gli appunti del corso compare un elenco di ciò che ``non viene mostrato'' dell'intelligenza artificiale: l'inquinamento delle miniere di litio, le terre rare estratte in zone di guerra spesso da manodopera minorile, le condizioni quasi schiavili della produzione elettronica, i data center alimentati da fonti non rinnovabili, i lavoratori di \emph{The Cleaners} costretti a visionare contenuti violenti. È quasi un riassunto dei primi due capitoli di Crawford, \emph{Earth} e \emph{Labor}, che ne offrono lo sviluppo narrativo e documentato in profondità.

\subsection{Echo chambers e social media come piazza pubblica}

Gli appunti discutono il caso Gonzalez contro Google, in cui i genitori di una ragazza morta in un attentato terroristico in Francia avevano fatto causa a Google sostenendo che l'algoritmo di personalizzazione di YouTube avesse contribuito alla radicalizzazione dei suoi assassini, e si chiedono se i social media siano una piazza pubblica. È lo stesso terreno su cui si muove il capitolo di Benanti sui social network e il mondo post-fattuale: l'economia dell'attenzione che premia indignazione e paura, le \emph{echo chamber} che isolano dal confronto, la perdita dell'intermediazione che mette sullo stesso piano un esperto e un complottista. Benanti offre qui un vocabolario teorico più articolato per un fenomeno che il corso introduce soprattutto attraverso casi giudiziari concreti.

\subsection{Pappagalli stocastici e la distinzione tra sintassi e semantica}

Gli appunti ricordano come Margaret Mitchell e Timnit Gebru siano state licenziate da Google poco prima dell'uscita di ChatGPT per aver messo in guardia sul rischio dei ``pappagalli stocastici'': i language model non sanno cosa dicono, ripetono parole che non comprendono. È esattamente l'argomento con cui Benanti distingue la sintassi delle macchine dalla semantica umana, e suggerisce che la critica filosofica di Benanti e l'avvertimento tecnico di Gebru e Mitchell, nato in un contesto aziendale molto concreto, condividano la stessa diagnosi di fondo.

\subsection{Tecnofeudalesimo e sovranità tecnologica}

Il corso descrive il tecnofeudalesimo osservando che Amazon non è più soltanto un mercato, ma un territorio su cui Jeff Bezos impone le proprie regole come un aristocratico che vive di rendita. È lo stesso fenomeno che Benanti descrive parlando di asimmetria di potere e di sovranità tecnologica: quando un algoritmo decide chi ha accesso a un servizio, sta esercitando una funzione che un tempo apparteneva solo allo stato. Crawford documenta empiricamente questa stessa erosione di sovranità nel capitolo sullo Stato; Benanti ne propone una lettura più esplicitamente politica e giuridica, fino a chiedere una governance pubblica capace di restituire alla politica il primato sull'efficienza tecnica.

\subsection{``La soluzione è solo tecnica, ma il problema è politico''}

Una delle osservazioni più nette degli appunti del corso, a proposito della censura adottata in modo diverso da Deepseek, ChatGPT e Grok, è che ``queste soluzioni sono solo tecniche, ma il problema è politico''. È precisamente il punto su cui Crawford e Benanti si dividono in modo più interessante: Benanti scommette sulla possibilità di tradurre principi etici in regole tecniche e in governance (l'algoretica), mentre Crawford diffida di questa traduzione e sposta l'attenzione sul potere economico e politico che resta intatto dietro qualsiasi soluzione tecnica. Discutere insieme questi due libri permette di mostrare che il problema posto dal corso non ha una sola risposta corretta, ma una tensione reale tra chi crede che le regole tecniche possano bastare e chi pensa che serva, prima di tutto, ridistribuire il potere.

%=============================================================================
\section{Tabella di sintesi}
%=============================================================================

\begin{table}[h!]
\centering
\renewcommand{\arraystretch}{1.3}
\begin{tabular}{p{4.2cm}p{5.3cm}p{5.3cm}}
\toprule
\textbf{Concetto del corso} & \textbf{In Crawford (Atlas of AI)} & \textbf{In Benanti (Le Macchine Sapienti)} \\
\midrule
Inevitabilismo tecnologico & ``Politics of refusal'' contro ``se può essere fatto, sarà fatto'' & L'IA come forma di potere che richiede scelte politiche, non un destino \\
Earth/Labor: cosa non si vede & Miniere di litio, cobalto, Mechanical Turk, ghost work & Lavoro e nuova competizione uomo-macchina nel processo produttivo \\
Echo chambers, Gonzalez vs Google & Capitolo Data: ``la fine del consenso'' & Economia dell'attenzione, echo chamber, perdita dell'intermediazione \\
Pappagalli stocastici, sintassi/semantica & ``L'IA non è intelligente'', solo inferenza statistica & Sintassi (macchina) vs semantica (uomo): la macchina non comprende \\
Tecnofeudalesimo, sovranità tecnologica & Capitolo State: NSA-ificazione, Palantir, potere privato & Asimmetria di potere, sovranità tecnologica, necessità di governance \\
Soluzione tecnica vs problema politico & Critica alle carte etiche, ``politics of refusal'' & Algoretica: codici etici e governance pubblica vincolante \\
Guerra autonoma e responsabilità & Project Maven, drone, ``uccidiamo sulla base di metadati'' & Artificial Moral Agent: chi è responsabile delle decisioni della macchina? \\
\bottomrule
\end{tabular}
\caption{Sintesi dei collegamenti tra il corso di Etica e i due libri.}
\end{table}

%=============================================================================
\section{Domande utili per l'esposizione orale}
%=============================================================================

Una traccia di domande con cui mettersi alla prova, utile per ripassare prima dell'esame:

\begin{enumerate}
  \item Spiega il senso della frase di Crawford ``l'intelligenza artificiale non è né artificiale né intelligente'' e confrontala con la distinzione di Benanti tra sintassi e semantica: in che senso sono la stessa critica, raccontata da due prospettive diverse?
  \item In che modo Crawford e Benanti usano entrambi il concetto di ``corpo'', ma per dire cose opposte? Spiega con il caso delle miniere di litio (Crawford) e con la riflessione sulla fragilità umana (Benanti).
  \item Spiega cos'è il ``passato che imprigiona il futuro'' secondo Benanti, e collega questo concetto alla genealogia che Crawford traccia tra Samuel Morton e ImageNet.
  \item Cos'è l'algoretica secondo Benanti? In che modo la sua proposta di governance si differenzia dalla ``politics of refusal'' di Crawford?
  \item Confronta il concetto di ``capitalismo di sorveglianza'' e ``sovranità tecnologica'' di Benanti con quello di ``registro di potere'' di Crawford. Che fenomeno descrivono entrambi?
  \item Usa il caso Gonzalez contro Google, discusso a lezione, per spiegare il capitolo di Benanti sui social network e il mondo post-fattuale.
  \item Perché, secondo gli appunti del corso, ``la soluzione è solo tecnica ma il problema è politico''? Mostra come questa frase si applichi sia al caso della censura nei diversi chatbot, sia al confronto tra l'approccio di Benanti e quello di Crawford.
\end{enumerate}

Un'ultima nota utile per l'orale: portare questo confronto tra i due libri durante l'esposizione mostra una capacità di sintesi che va oltre il semplice riassunto di un singolo testo, ed è proprio il tipo di collegamento autonomo che un esame orale valorizza.

\end{document}
